% Section 5.2: Size of the algebra and relation to $\bigwedge V$ | Part: 1 | Depends on: sec:metric, sec:ext | Status: done | Original lines: 1317-1344
%!TEX root = ../../../main.tex
\subsection{Size of the algebra and relation to $\bigwedge V$}

\begin{theorem}[Dimension, basis]\label{thm:cliffbasis}
$\dim\Cl(V,Q)=2^n$, with basis $e_I=e_{i_1}e_{i_2}\cdots e_{i_k}$ for $I=\{i_1<\cdots<i_k\}\subseteq\{1,\dots,n\}$ (and $e_\emptyset=1$).
\end{theorem}
\begin{proof}[Proof sketch]
\emph{Upper bound.} Using $e_ie_j=-e_je_i$ and $e_i^2=\varepsilon_i$, any word can be reordered into $\pm$ (a scalar) $\cdot\,e_I$ for some $I$, so the $e_I$ span.
\emph{Lower bound (Chevalley's construction).} For $v\in V$ let $\iota_v$ be the \emph{interior product} on $\bigwedge V$: the unique antiderivation of degree $-1$ with $\iota_v(w)=g(v,w)$ for $w\in V$. Define
$\rho(v)\,\omega=v\wedge\omega+\iota_v\omega$ on $\bigwedge V$. Using $\iota_v(v\wedge\omega)=g(v,v)\omega-v\wedge\iota_v\omega$ and $\iota_v^2=0$, $v\wedge v=0$ one finds
\[
 \rho(v)^2\omega=v\wedge\iota_v\omega+\iota_v(v\wedge\omega)=Q(v)\,\omega .
\]
By Theorem~\ref{thm:cliffuniv}, $\rho$ extends to an algebra map $\Cl(V,Q)\to\mathrm{End}(\bigwedge V)$. The map $x\mapsto\rho(x)1$ sends $e_I\mapsto e_{i_1}\wedge\cdots\wedge e_{i_k}$ (for an orthogonal basis the interior terms vanish), which are linearly independent by Theorem~\ref{thm:extbasis}. Hence the $e_I$ are independent.
\end{proof}

\begin{corollary}[Grades]\label{cor:grades}
The map $\bigwedge V\to\Cl(V,Q)$ given by antisymmetrization,
\[
 v_1\wedge\cdots\wedge v_k\;\longmapsto\;\frac1{k!}\sum_{\sigma\in S_k}\sgn(\sigma)\,v_{\sigma(1)}v_{\sigma(2)}\cdots v_{\sigma(k)},
\]
is a (basis-independent) linear isomorphism of vector spaces. We use it to write $\GA(V)=\bigoplus_{k=0}^n\GA^k(V)$ and to define the
\emph{grade projections} $\gr{A}{k}$. For \emph{mutually orthogonal} vectors, $a_1a_2\cdots a_k=a_1\wedge\cdots\wedge a_k$. The product does \emph{not} preserve grade, but it preserves the parity (even/odd), so the \emph{even part} $\GA^+=\bigoplus_{k\ \mathrm{even}}\GA^k$ is a subalgebra.
\end{corollary}

\begin{remark}
So, as vector spaces, $\Cl(V,Q)\cong\bigwedge V$; they differ only in the \emph{multiplication}. When $Q=0$ the Clifford algebra is the exterior algebra. In this sense the Clifford algebra is a \emph{deformation} of $\bigwedge V$ by the metric.
\end{remark}
